% !TeX root = ../../Book2.tex
% 纲要标题：### 15.3 投射覆盖
% 纲要位置：计划纲要.md，第 1541 行。

\section{投射覆盖}
\label{b2:sec:1503}

自由满射总存在，但其来源可能带有可以删去的投射分量。投射覆盖要求满射在模论意义下已经不能缩小来源。为了构造它，先要把半单商中的列模提升回来；这一步依赖幂等元提升，不能用商中的分解代替原代数中的分裂。本节的 \(A\) 仍有限维，\(J=J(A)\)；幂等元提升的两个定理则只需所写的环与幂零理想假设。

% 纲要要点（供目录核对）：
%
% * 极小生成与本质满射
% * 在构造投射覆盖前，就地证明幂零理想模的幂等元提升，并处理有限组正交幂等元的相容提升
% * 明确本质满射的模论含义为核是多余子模；从提升后的本原幂等元构造不可分解投射模及其到简单模的覆盖
% * 有限维代数中的投射覆盖；目标有限生成时来源的有限性，投射覆盖与最少自由生成的区别
% * 不可分解投射模与简单模
% * 由有限维性、半单商及根的幂零性逐项验证覆盖的存在和唯一性，避免把半单商中的分解直接当作原模中的分裂
% * 正文给出分别提升不保持正交的反例；唯一性包含目标同构时每个提升可逆，极小性说明投射目标的覆盖为同构
%

\subsection{极小生成与本质满射}
\label{b2:subsec:1503:01}

\begin{definition}\label{b2:def:superfluous-submodule}
设 \(R\) 为含幺环、\(P\) 为幺左 \(R\) 模、\(N\le P\) 为子模。若对每个子模 \(L\le P\)，等式 \(L+N=P\) 都蕴含 \(L=P\)，则称 \(N\) 为 \(P\) 的\term[duoyu-zimo]{多余子模}[superfluous submodule]，记为 \(N\ll P\)。若幺左 \(R\) 模的满同态 \(p:P\to M\) 满足 \(\ker p\ll P\)，则称 \(p\) 为\term[benzhi-manshe]{本质满射}[essential epimorphism]。
\end{definition}
\symidx{\(N\ll P\)：\(N\) 是 \(P\) 的多余子模}

这一条件也可写成：对每个子模 \(L\le P\)，若 \(p(L)=M\)，则 \(L=P\)。事实上，\(p(L)=M\) 等价于每个 \(x\in P\) 可写成 \(\ell+n\)，其中 \(\ell\in L,n\in\ker p\)，也就是 \(L+\ker p=P\)。因此任何仍能满射到目标的子模都已经是整个来源。多余性总是相对于指定的母模而言，判定时量化的是这个母模的全部子模。

\ProseParagraphBreak

等价地，对任意左 \(R\) 模同态 \(g:Q\to P\)，若复合 \(pg:Q\to M\) 满射，则 \(g\) 满射：此时 \(p(g(Q))=M\)，上一条件给出 \(g(Q)=P\)。反过来，把 \(g\) 取为任意子模 \(L\le P\) 的包含映射，就得到上一条件。

\begin{definition}\label{b2:def:projective-cover}
设 \(R\) 为含幺环、\(M\) 为幺左 \(R\) 模。若幺左 \(R\) 模的满同态 \(p:P\to M\) 的来源 \(P\) 投射，且 \(\ker p\ll P\)，则称 \(p\) 为 \(M\) 的一个\term[toushe-fugai]{投射覆盖}[projective cover]。定义中保留这个满同态，而不仅记录 \(P\) 的同构类型。
\end{definition}

若 \(M\) 有限生成，它的任意投射覆盖的来源也有限生成：取 \(M\) 的有限生成元，逐个提升到 \(P\)，令 \(L\) 为这些提升所生成的子模，则 \(p(L)=M\)。本质满射的条件给出 \(L=P\)。因此，在有限维代数 \(A\) 上，有限生成目标模的覆盖来源有限生成，进而有限维；这适用于本章讨论的有限维目标模。

\ProseParagraphBreak
\cref{b2:cor:nakayama-generators}说明，\(M/JM\) 的一组 \(A/J\) 模生成元可以提升为 \(M\) 的生成元。局部代数中，剩余模是在除环 \(A/J\) 上的向量空间，选基可得到最少的生成元。但一般情形下，由最少生成元给出的自由满射仍未必是投射覆盖。

\ProseParagraphBreak
例如，设 \(k\) 为域、\(A=M_2(k)\)、\(S=k^2\)。\(S\) 由一个非零向量生成，取第一列的映射 \(A\to S\) 因而用了最少的一个模生成元。不过真左理想 \(AE_{11}\) 已经满射到 \(S\)，故这个自由满射不是本质满射。限制 \(AE_{11}\to S\) 则是同构；\(AE_{11}\) 是正则模的直和因子，因而投射，这才给出一个投射覆盖。\(AE_{11}\) 的 \(k\) 维数为二，不能是 \(A\) 上的自由模。因此，最少生成元使自由来源的秩最小，来源仍可能缩小为一个非自由的投射模。

\subsection{幂等元与正交幂等元组的提升}
\label{b2:subsec:1503:02}

\begin{theorem}\label{b2:thm:idempotent-lifting}
设 \(R\) 为含幺环，\(I\) 为幂零双边理想。每个幂等元 \(\bar e\in R/I\) 都能提升为幂等元 \(e\in R\)。此外，两个这样的提升 \(e,f\) 共轭：存在 \(u\in1+I\)，使 \(f=ueu^{-1}\)。幂零假设不能一般省去：在 \(R=\mathbb Z\)、\(I=6\mathbb Z\) 中，幂等元 \(3+6\mathbb Z\) 不能提升为整数幂等元。
\end{theorem}
\begin{proof}
先处理 \(I^2=0\)。任取 \(a\) 提升 \(\bar e\)，令 \(d=a^2-a\in I\)，这是 \(a\) 未必幂等的误差。希望取 \(e=a+c\)，其中 \(c\in I\) 且与 \(a\) 交换。因为 \(I^2=0\) 使 \(c^2=0\)，所需条件化为
\[
e^2-e=d+ac+ca-c=d+(2a-1)c=0.
\]
这个误差修正方程对 \(c\) 是线性的。又有 \((2a-1)^2=1+4d\)；对任意 \(x\in I\)，\(dx\in I^2=0\)，所以 \((2a-1)^2x=x\)。因此在 \(I\) 上左乘 \(2a-1\) 的算子平方为恒等，可以取
\[
c=-(2a-1)d=(1-2a)d.
\]
所取 \(d,c\) 都是 \(a\) 的多项式，故确实与 \(a\) 交换；它们也都在 \(I\) 中。代入得到
\[
e^2-e=d+(2a-1)c
=d-(2a-1)^2d.
\]
由于 \((2a-1)^2=1+4d\)，右侧等于 \(-4d^2=0\)。因此 \(e\) 幂等，且与 \(a\) 模 \(I\) 相同。这里没有除以二，对任意特征都成立。

一般地，取 \(I^N=0\)。已在 \(R/I^r\) 中得到的幂等元，可以提升到 \(R/I^{r+1}\)：自然满射的核为 \(I^r/I^{r+1}\)，它的平方为零，因为 \(2r\ge r+1\)、\(r\ge1\)。将平方零情形用于这个商环，依次作 \(r=1,\ldots,N-1\) 的提升，最后得到 \(R/I^N=R\) 中的幂等元。

为证共轭性，先把 \(e,f\) 看成正则右模 \(R_R\) 上的左乘投影；左乘是右 \(R\) 线性的，因为 \((ex)r=e(xr)\)。它们分别给出分解
\[
R_R=eR\oplus(1-e)R=fR\oplus(1-f)R.
\]
我们寻找的左乘映射应把第一组分解的两个部分分别送到第二组对应的部分，从而满足 \(ue=fu\)。对 \(eR\) 中的元素左乘 \(f\)，对 \((1-e)R\) 中的元素左乘 \(1-f\)，便得到候选元素
\[
u=fe+(1-f)(1-e).
\]
事实上，对 \(x\in eR\) 有 \(ux=fx\in fR\)，对 \(x\in(1-e)R\) 有 \(ux=(1-f)x\in(1-f)R\)；直接相乘也得到 \(ue=fe=fu\)。这个公式先完成分解的匹配，还须保证匹配映射可逆。因 \(\bar f=\bar e\)，模 \(I\) 后有 \(\bar u=\bar e+(1-\bar e)=1\)，故 \(u\in1+I\)。幂零性使 \(u\) 可逆，可由有限几何和写出其逆。于是 \(ue=fu\) 给出 \(f=ueu^{-1}\)。

最后，\(3^2-3=6\)，所以 \(3+6\mathbb Z\) 幂等。整数幂等元满足 \(e(e-1)=0\)，而 \(\mathbb Z\) 是整环，故只有零或一；它们模六都不等于三，说明这个幂等元不能提升。这里 \((6\mathbb Z)^r=6^r\mathbb Z\ne0\) 对每个正整数 \(r\) 成立，恰好缺少幂零假设。
\end{proof}

要提升一组分解，所选提升还须保持彼此正交和总和为一；分别任意提升每个幂等元并不足够。\cref{b2:prop:idempotent-lift-compatibility-counterexample}在本节末给出各自幂等而不相容的完整计算。下一定理把这种相容性一并保证。

\begin{theorem}\label{b2:thm:orthogonal-idempotent-lifting}
设 \(R\) 为含幺环、\(I\subseteq R\) 为幂零双边理想，\(m\) 为正整数。若 \(\bar e_1,\ldots,\bar e_m\in R/I\) 满足
\[
\bar e_i\bar e_j=\delta_{ij}\bar e_i,
\qquad \sum_{i=1}^m\bar e_i=1,
\]
则存在提升 \(e_1,\ldots,e_m\in R\)，同样满足 \(e_ie_j=\delta_{ij}e_i\) 与 \(\sum_i e_i=1\)。
\end{theorem}
\begin{proof}
对 \(m\) 归纳。\(m=1\) 时取 \(e_1=1\)。设 \(m>1\)，先用\cref{b2:thm:idempotent-lifting}提升 \(\bar e_1\)，并令 \(c=1-e_1\)。余下部分的单位应为 \(c\)，所以转到角环 \(cRc\) 中提升其余幂等元。映射
\[
cRc\longrightarrow\bar c(R/I)\bar c
\]
满射：右侧元素可由任意代表元两边乘 \(c\) 得到。其核为 \(cIc\)，因为同时属于 \(I\) 与 \(cRc\) 的元素 \(x\) 满足 \(x=cxc\)。这个核幂零，\((cIc)^N\subseteq cI^Nc=0\)。

余下的 \(\bar e_2,\ldots,\bar e_m\) 都在该商角环中，仍相互正交，总和为 \(\bar c\)。在 \(cRc\) 中应用归纳假设，得到正交提升 \(e_2,\ldots,e_m\)，总和为 \(c\)。它们均被 \(e_1\) 在左右两侧乘为零，所以连同 \(e_1\) 得到所求完整正交组。
\end{proof}

平方零修正与角环归纳可参阅 Etingof 等人的《Introduction to Representation Theory》\cite[Proposition 9.1.1; Corollary 9.1.3, pp. 209--210]{EtingofEtAl2011RepresentationTheory}。本书已把单个提升和整组相容提升分别展开；后面的构造将同时使用这两个结论。

\subsection{多余子模与简单模的投射覆盖}
\label{b2:subsec:1503:03}

\begin{proposition}\label{b2:prop:finite-superfluous-radical}
设 \(k\) 为域、\(A\) 为有限维结合含幺 \(k\) 代数，\(J=J(A)\)。有限维幺左 \(A\) 模 \(P\) 的子模 \(N\) 多余，当且仅当 \(N\subseteq JP=\operatorname{rad}P\)。因此，对有限生成投射幺左 \(A\) 模的满射 \(p:P\to M\)，以下条件等价：它是投射覆盖；\(\ker p\subseteq JP\)；诱导映射 \(P/JP\to M/JM\) 是同构。若这个满射是投射覆盖且目标 \(M\) 本身投射，则 \(p\) 为同构。
\end{proposition}
\begin{proof}
若 \(N\ll P\)，而某个极大子模 \(L\) 不包含 \(N\)，则 \(L\subsetneq L+N\)，极大性迫使 \(L+N=P\)，与 \(L\ne P\) 矛盾。因此 \(N\) 包含于所有极大子模的交，由\cref{b2:thm:module-radical-jm}得到 \(N\subseteq JP=\operatorname{rad}P\)。

反过来，若 \(N\subseteq JP\) 且 \(L+N=P\)，则 \(L+JP=P\)。商 \(P/L\) 有限生成，并满足 \(J(P/L)=P/L\)，由\cref{b2:thm:noncommutative-nakayama}，\(P/L=0\)，即 \(L=P\)。

对满射 \(p\)，有 \(p(JP)=JM\)。诱导的商映射满射，核为 \((\ker p+JP)/JP\)：若 \(p(x)\in JM\)，取 \(y\in JP\) 使 \(p(y)=p(x)\)，便有 \(x-y\in\ker p\)。所以该核为零恰等于 \(\ker p\subseteq JP\)，得到其余等价性。

若 \(p\) 是覆盖且 \(M\) 投射，取截面 \(s:M\to P\)，便有 \(P=\ker p\oplus s(M)\)。核多余使 \(s(M)=P\)，从而 \(\ker p=0\)，\(p\) 是同构。
\end{proof}

\begin{definition}\label{b2:def:primitive-idempotent}
设 \(k\) 为域、\(A\) 为结合含幺 \(k\) 代数、\(e\in A\) 为非零幂等元。若 \(e\) 不能写成两个非零正交幂等元之和，则称它为\term[benyuan-midengyuan]{本原幂等元}[primitive idempotent]。这里的正交要求两个方向的乘积都为零。
\end{definition}

这里“本原”修饰幂等元，表示它不能再作正交分解；\cref{b2:def:primitive-semiprimitive}中的本原环则要求环有忠实简单模，两者的定义不同。

\begin{proposition}\label{b2:prop:primitive-projective-indecomposable}
设 \(k\) 为域、\(A\) 为结合含幺 \(k\) 代数、\(e\in A\) 为非零幂等元。幺左 \(A\) 模 \(Ae\) 有限生成且投射；它不可分解，当且仅当 \(e\) 本原。
\end{proposition}
\begin{proof}
\(Ae\) 由单个元素 \(e\) 生成。正则左模具有分解 \(A=Ae\oplus A(1-e)\)，投影为右乘 \(e\)，所以 \(Ae\) 是有限自由模的直和因子，因而投射。若 \(e=f+g\) 是非零正交幂等元之和，则 \(Ae=Af\oplus Ag\)，给出非平凡分解。

反过来，若 \(Ae\) 有非平凡分解，对应投影是 \(\operatorname{End}_A(Ae)\) 中一个非零且非恒等的幂等元。由\cref{b2:prop:corner-endomorphism}，该自同态环同构于 \((eAe)^{\mathrm{op}}\)，所以角环中也有 \(f^2=f\)、\(f\ne0,e\)。于是 \(e=f+(e-f)\)，且 \(f(e-f)=(e-f)f=0\)，这与本原性矛盾。
\end{proof}

\begin{proposition}\label{b2:prop:primitive-projective-cover}
设 \(k\) 为域、\(A\) 为有限维结合含幺 \(k\) 代数，\(J=J(A)\)。取 \(A/J\) 的本原幂等元 \(\bar e\)，并取其在 \(A\) 中的幂等元提升 \(e\)，则 \(e\) 本原，\(Ae/Je\) 是简单幺左 \(A\) 模，且自然满射
\[
Ae\longrightarrow Ae/Je
\]
是投射覆盖。每个简单幺左 \(A\) 模都能由这样的满射得到覆盖。
\end{proposition}
\begin{proof}
先看单个矩阵环 \(M_n(D)\)，其中 \(D\) 为除环。幂等矩阵作为右 \(D\) 向量空间上的算子，其像与核给出直和；选择相应基后，它共轭于 \(\operatorname{diag}(I_r,0)\)。后者是 \(r\) 个秩一坐标投影之和，所以 \(r>1\) 时不本原。若 \(r=1\)，两个非零正交幂等元的像会给出两个非零直和分量，不能同时放在这个一维像中。因此非零幂等元本原当且仅当 \(r=1\)。再看\eqref{b2:eq:finite-algebra-semisimple-quotient}中的有限矩阵环直积，本原幂等元只能在一个矩阵因子中具有秩一，其余因子为零。相应左理想同构于一个列模，且简单：若列 \(v\) 的第 \(j\) 个坐标非零，为把它映到任意给定列 \(w\)，取矩阵 \(C\) 的第 \(j\) 列为 \((w_iv_j^{-1})_i\)，其余列为零，便有 \(Cv=w\)。

\(J\) 零化每个简单左 \(A\) 模，所以它们都可看作 \(A/J\) 的简单模。半单商的正则模是上述简单列模的有限直和；任意简单模都是这个正则模的商。相应满射在至少一个列模上的限制非零，由\cref{b2:lem:schur}，该限制是同构。这也说明上述构造遍及全部简单模。

若 \(e=f+g\) 是正交幂等元分解，模 \(J\) 后，\(\bar e\) 的本原性使 \(\bar f,\bar g\) 至少一个为零。该幂等元便属于幂零理想 \(J\)，只能为零，因为幂等元的每个正幂都等于自身。因此 \(e\) 本原。最后，\(Ae\to(A/J)\bar e\) 的核为 \(Je\)：核元素既在 \(J\) 中，又满足 \(x=xe\)，故在 \(Je\) 中。于是 \(Ae/Je\) 是刚才的简单列模，且核为 \(J(Ae)\)。由\cref{b2:prop:finite-superfluous-radical}，满射为投射覆盖。
\end{proof}

\subsection{投射覆盖的存在性}
\label{b2:subsec:1503:04}

从 \(A/J\) 的每个矩阵因子选出全部对角矩阵单位，利用\cref{b2:thm:orthogonal-idempotent-lifting}得到 \(A\) 中的完整正交提升。每个提升给出上述简单列模的覆盖；同一矩阵因子的列模同构，但暂不预设其提升模已经同构。

\ProseParagraphBreak
从每一种简单模的同构类型中取一个代表 \(S_1,\ldots,S_s\)，并固定由\cref{b2:prop:primitive-projective-cover}构造的覆盖 \(p_i:P_i\to S_i\)。于是 \(P_i/JP_i\cong S_i\)，且 \(P_i\) 不可分解。

\begin{theorem}\label{b2:thm:projective-cover-existence}
设 \(k\) 为域、\(A\) 为有限维结合含幺 \(k\) 代数，\(J=J(A)\)。从全部简单幺左 \(A\) 模的同构类中各选一个代表 \(S_1,\ldots,S_s\)，并给每个 \(S_i\) 选定投射覆盖 \(p_i:P_i\to S_i\)，则每个有限维幺左 \(A\) 模 \(M\) 都有投射覆盖。若
\[
M/JM\cong\bigoplus_{i=1}^s S_i^{\oplus m_i},
\]
可以取覆盖的来源为 \(P=\bigoplus_iP_i^{\oplus m_i}\)。
\end{theorem}
\begin{proof}
由\cref{b2:thm:module-radical-jm}，所写的有限半单分解存在。合并 \(p_i\) 的相应份数并接上这个同构，得到满同态 \(h:P\to M/JM\)，其在商上诱导同构 \(P/JP\cong M/JM\)。\(P\) 是有限个投射模的直和，因而投射，可以沿商映射 \(\pi:M\to M/JM\) 提升 \(h\)，得到 \(p:P\to M\) 且 \(\pi p=h\)。

因为 \(h\) 满射，有 \(p(P)+JM=M\)。对有限生成的 \(M\) 及其子模 \(p(P)\) 应用\cref{b2:cor:nakayama-generators}的生成元提升形式，得到 \(p(P)=M\)。又因为 \(p\) 在模 \(J\) 后是同构，\(\ker p\subseteq JP\)。由\cref{b2:prop:finite-superfluous-radical}，这个满射正是投射覆盖。
\end{proof}

这里并未把 \(M/JM\) 的简单直和分解直接提升成 \(M\) 的简单直和分解。被提升的是它们的投射覆盖，以及一个从投射模出发的映射；满射性由 Nakayama 引理重新保证。

\ProseParagraphBreak
例如在 \(A=k[t]/(t^3)\) 中，\(M=A/(\bar t^2)\) 的最大半单商为 \(k\)，对应投射覆盖是 \(A\to M\)。核 \((\bar t^2)\) 包含在 \(JA\) 中，所以多余；这个覆盖保留了 \(M\) 的两层结构，而来源 \(A\) 本身有三层。

\subsection{不可分解投射模与简单模}
\label{b2:subsec:1503:05}

\begin{theorem}\label{b2:thm:indecomposable-projective-simple}
设 \(k\) 为域、\(A\) 为有限维结合含幺 \(k\) 代数，\(J=J(A)\)。从全部简单幺左 \(A\) 模的同构类中各选一个代表 \(S_1,\ldots,S_s\)，并按\cref{b2:prop:primitive-projective-cover}给每个 \(S_i\) 选定形如 \(p_i:P_i=Ae_i\to S_i\) 的投射覆盖。于是每个有限生成投射幺左 \(A\) 模都是 \(P_1,\ldots,P_s\) 的有限直和。有限生成不可分解投射模恰为这些 \(P_i\) 的同构类型，而
\[
P\longmapsto P/JP
\]
给出有限生成不可分解投射模与简单模的同构类型之间的一一对应。
\end{theorem}
\begin{proof}
对有限生成投射模 \(Q\)，按\cref{b2:thm:projective-cover-existence}构造覆盖 \(p:P\to Q\)，其中 \(P\) 是若干 \(P_i\) 的直和。\(Q\) 投射，而 \(p\) 是投射覆盖，由\cref{b2:prop:finite-superfluous-radical}，\(p\) 本身就是同构，故 \(Q\cong P\)。

各 \(P_i\) 已由\cref{b2:prop:primitive-projective-cover}证明不可分解；若 \(Q\) 不可分解，上述直和只能有一个非零分量。不同的 \(P_i\) 不同构，因为其商 \(P_i/JP_i\cong S_i\) 不同构。每个简单模都有其中一个覆盖，故得到双射。
\end{proof}

因此，这里对应的是三种对象的同构类型：
\[
S_i\quad\longleftrightarrow\quad P_i\cong Ae_i
\quad\longleftrightarrow\quad
\{\,e\text{ 为本原幂等元}\mid Ae\cong Ae_i\,\}.
\]
右端是一类本原幂等元，而不是某一个唯一的幂等元。

\ProseParagraphBreak
直和中的重数也可以直接从最大半单商读出。对有限生成投射左 \(A\) 模 \(Q\)，令 \(\Delta_i=\operatorname{End}_A(S_i)\)，它由\cref{b2:lem:schur}为除环。\(\operatorname{Hom}_A(Q,S_i)\) 通过像端的复合成为左 \(\Delta_i\) 向量空间，具体为 \(\alpha\cdot f=\alpha\circ f\)。因为 \(J\) 零化 \(S_i\)，每个同态都通过 \(Q/JQ\) 分解，因而
\begin{equation}\label{b2:eq:projective-simple-multiplicity}
Q\cong\bigoplus_iP_i^{\oplus m_i}
\quad\Longrightarrow\quad
m_i=\dim_{\Delta_i}\operatorname{Hom}_A(Q,S_i).
\end{equation}
理由是不同简单模之间没有非零同态，\(\operatorname{Hom}_A(S_i,S_i)=\Delta_i\) 在自身上的左维数为一。这也证明这些重数唯一，且没有假设 \(\Delta_i=k\)。\cref{b2:prop:semisimple-multiplicity}从来源简单模的自同态作预复合，得到右除环向量空间；这里简单模位于像端，作后复合，因而得到左除环向量空间。

\subsection{投射覆盖的唯一性与极小性}
\label{b2:subsec:1503:06}

存在性已经由\cref{b2:thm:projective-cover-existence}得到。以下唯一性须比较连同映射的覆盖；只比较来源的维数，不足以说明它们在 \(M\) 上相容。

\begin{theorem}\label{b2:thm:projective-cover-uniqueness}
设 \(k\) 为域、\(A\) 为有限维结合含幺 \(k\) 代数，\(M,N\) 为有限维幺左 \(A\) 模，\(h:M\to N\) 为 \(A\) 模同构，\(p:P\to M\) 与 \(q:Q\to N\) 为投射覆盖。则存在满足 \(qa=hp\) 的 \(A\) 模同态 \(a:P\to Q\)，且每个满足这一等式的 \(a\) 都是同构。取 \(N=M\)、\(h=1_M\)，便得到两个覆盖在目标 \(M\) 上同构；这样的同构本身未必唯一。
\end{theorem}
\begin{proof}
由 \(P\) 的投射性，可以沿满射 \(q\) 提升 \(hp\)，所以满足 \(qa=hp\) 的 \(a\) 存在。现在固定任意一个这样的 \(a\)。由 \(Q\) 的投射性，沿 \(p\) 提升 \(h^{-1}q\)，得到 \(b:Q\to P\)，满足 \(pb=h^{-1}q\)。于是 \(pba=p\)、\(qab=q\)。对任意 \(x\in P\)，有 \(x-ba(x)\in\ker p\)，所以
\[
ba(P)+\ker p=P.
\]
核多余，故 \(ba(P)=P\)。同样由 \(qab=q\) 及 \(\ker q\ll Q\) 得 \(ab(Q)=Q\)。这两步只用了覆盖的多余核条件。

目标 \(M,N\) 有限生成，本质满射性使来源 \(P,Q\) 也有限生成；又因 \(A\) 有限维，\(P,Q\) 都有限维，因而具有有限长度。\cref{b2:prop:finite-length-endomorphism}使满自同态 \(ba\) 与 \(ab\) 可逆。于是 \((ba)^{-1}b\) 是 \(a\) 的左逆，而 \(b(ab)^{-1}\) 是其右逆；一个映射的左逆与右逆相同，故给定的 \(a\) 是同构。

例如 \(A=k[t]/(t^2)\) 的覆盖 \(A\to k\) 上，恒等映射和乘以 \(1+\bar t\) 的映射都是同构，且模 \(\bar t\) 后都诱导恒等，二者不同。因此唯一的是覆盖在目标上同构的类型。
\end{proof}

\begin{proposition}\label{b2:prop:projective-cover-minimality}
设 \(k\) 为域、\(A\) 为有限维结合含幺 \(k\) 代数，\(M\) 为幺左 \(A\) 模。对有限生成投射幺左 \(A\) 模的满射 \(p:P\to M\)，它是投射覆盖，当且仅当每个满足 \(pu=p\) 的 \(A\) 模自同态 \(u:P\to P\) 都可逆。若 \(p\) 是覆盖，且 \(q:Q\to M\) 是任意投射幺左 \(A\) 模的满射，则 \(P\) 是 \(Q\) 的直和因子，不要求 \(Q\) 有限生成。
\end{proposition}
\begin{proof}
覆盖情形中，\(u(P)+\ker p=P\)，所以 \(u\) 满射。由命题假设，\(P\) 在有限维 \(k\) 代数 \(A\) 上有限生成，故它本身有限维，\(u\) 因而可逆。反过来，若 \(L+\ker p=P\)，则 \(p|_L:L\to M\) 满射。利用 \(P\) 的投射性，把 \(p\) 提升到 \(L\)，再复合包含映射得到 \(u:P\to P\)，其像在 \(L\) 中且 \(pu=p\)。若它必可逆，则 \(L=P\)，所以核多余。

对第二个断言，投射性给出 \(a:P\to Q\)、\(b:Q\to P\)，满足 \(qa=p\)、\(pb=q\)。于是 \(pba=p\)，第一部分使 \(ba\) 可逆。映射 \(a(ba)^{-1}:P\to Q\) 是 \(b\) 的截面，因此 \(Q\cong P\oplus\ker b\)。
\end{proof}

这使“覆盖极小”有了可检验的含义：任何其他投射满射的来源都含有它作为直和因子；若一个来源还能通过自同态缩小却仍覆盖目标，就不满足这种极小性。半单代数上的有限维模都投射，所以它们的覆盖不会增添额外的层，来源经覆盖映射与目标同构。关于同构类的选择虽然唯一，关于具体提升映射的选择仍须在后续构造中保留。

\begin{proposition}\label{b2:prop:idempotent-lift-compatibility-counterexample}
设 \(k\) 为域、\(R=M_2\bigl(k[\varepsilon]/(\varepsilon^2)\bigr)\)，\(\varepsilon\) 的商类仍记为 \(\varepsilon\)，\(E_{ij}\) 为矩阵单位，\(I=\varepsilon R\)。虽然 \(I^2=0\)，幂等元
\[
e=E_{11}+\varepsilon E_{12},\qquad f=E_{22}
\]
分别提升了 \(R/I\) 中正交且和为一的 \(\bar E_{11},\bar E_{22}\)，所选 \(e,f\) 却不正交。若改取 \(f'=1-e=E_{22}-\varepsilon E_{12}\)，则 \(e,f'\) 为同一商幂等元组的正交提升，且和为一。
\end{proposition}
\begin{proof}
\(I\) 为双边理想，且 \(\varepsilon^2=0\) 给出 \(I^2=0\)。由矩阵单位乘法，\(E_{11}E_{12}=E_{12}\)、\(E_{12}E_{11}=0\)，所以
\[
e^2=E_{11}+\varepsilon E_{12}=e,\qquad f^2=f.
\]
模 \(I\) 后，它们分别为 \(E_{11},E_{22}\)，其乘积在两个方向都为零，和为单位。但是在 \(R\) 中，
\[
ef=\varepsilon E_{12}\ne0,\qquad fe=0.
\]
因此分别取出的幂等提升不相互正交，它们的和也为 \(1+\varepsilon E_{12}\)，而不是一。取 \(f'=1-e\)，由 \(e^2=e\) 直接得到 \((f')^2=f'\)、\(ef'=f'e=0\) 及 \(e+f'=1\)；同时 \(f'-f\in I\)，所以没有改变第二个商幂等元。所有计算均适用于任意特征。
\end{proof}
